% Standalone problem document, generated from the adjacent README.md.
% Compile with XeLaTeX. Mathematical content is maintained in Markdown.
\documentclass[11pt,a4paper]{article}
\usepackage[fontsize=10.5pt]{scrextend}
\usepackage[margin=22mm,headheight=15pt,headsep=9mm,footskip=11mm]{geometry}
\usepackage{fontspec}
\setmainfont{texgyrepagella}[Extension=.otf,UprightFont=*-regular,BoldFont=*-bold,ItalicFont=*-italic,BoldItalicFont=*-bolditalic]
\setsansfont{texgyreheros}[Extension=.otf,UprightFont=*-regular,BoldFont=*-bold,ItalicFont=*-italic,BoldItalicFont=*-bolditalic]
\setmonofont{texgyrecursor}[Extension=.otf,UprightFont=*-regular,BoldFont=*-bold,ItalicFont=*-italic,BoldItalicFont=*-bolditalic,Scale=MatchLowercase]
\usepackage{amsmath,amssymb}
\usepackage{unicode-math}
\setmathfont{texgyrepagella-math.otf}
\usepackage{xcolor,microtype,enumitem,needspace,fancyhdr}
\usepackage{bookmark}
\definecolor{ink}{HTML}{17344B}
\definecolor{accent}{HTML}{197D80}
\definecolor{muted}{HTML}{596773}
\hypersetup{colorlinks=true,linkcolor=accent,urlcolor=accent,pdftitle={FR-02 - Sharp
sample complexity for restricted isometries from the cyclic Fourier
matrix},pdfauthor={Open Problems in Numerical Linear Algebra}}
\urlstyle{same}
\setlength{\parindent}{0pt}
\setlength{\parskip}{5pt plus 1pt minus 1pt}
\linespread{1.04}
\setlength{\emergencystretch}{3em}
\widowpenalty=10000
\clubpenalty=10000
\displaywidowpenalty=10000
\allowdisplaybreaks[1]
\setlist{nosep,leftmargin=1.5em}
\providecommand{\tightlist}{\setlength{\itemsep}{0pt}\setlength{\parskip}{0pt}}
\providecommand{\pandocbounded}[1]{#1}
\pagestyle{fancy}
\fancyhf{}
\fancyhead[L]{\sffamily\footnotesize\color{muted}OPEN PROBLEMS IN NUMERICAL LINEAR ALGEBRA}
\fancyhead[R]{\sffamily\bfseries\footnotesize\color{accent}FR-02}
\fancyfoot[L]{\parbox[b]{0.9\textwidth}{\sffamily\fontsize{8}{10}\selectfont\color{muted}Editorial ratings; a bounded search cannot certify that a problem remains unsolved.\\Literature check: 2026-09-08}}
\fancyfoot[R]{\sffamily\footnotesize\color{muted}\thepage}
\renewcommand{\headrulewidth}{0.3pt}
\renewcommand{\headrule}{\hbox to\headwidth{\color{accent}\leaders\hrule height \headrulewidth\hfill}}
\makeatletter
\renewcommand{\section}{\Needspace{5\baselineskip}\@startsection{section}{1}{0pt}{12pt plus 2pt minus 2pt}{4pt}{\sffamily\bfseries\large\color{ink}}}
\renewcommand{\subsection}{\Needspace{4\baselineskip}\@startsection{subsection}{2}{0pt}{9pt plus 1pt minus 1pt}{3pt}{\sffamily\bfseries\normalsize\color{ink}}}
\makeatother
\setcounter{secnumdepth}{0}
\begin{document}
{\sffamily\small\color{muted}Frames and matrix designs\par}
\vspace{3pt}
{\raggedright\hyphenpenalty=10000\exhyphenpenalty=10000\sffamily\bfseries\fontsize{21}{24}\selectfont\color{ink}Sharp
sample complexity for restricted isometries from the cyclic Fourier
matrix\par}
\vspace{7pt}
{\sffamily\small\textbf{Difficulty:} extreme\quad\textbf{Importance:} broadly
interesting\par}
\vspace{4pt}
{\color{accent}\hrule height 0.8pt}
\vspace{6pt}
\textbf{Status:} source-stated sharp-order open problem; no resolution
found in screening on 2026-09-08.

Let \(F_N\in\mathbb C^{N\times N}\) be the unitary cyclic discrete
Fourier matrix, \[
(F_N)_{j\ell}=N^{-1/2}\exp(-2\pi\mathrm i j\ell/N),\quad0\leq j,\ell<N.
\] Choose \(r_1,\ldots,r_m\) independently and uniformly from
\(\{0,\ldots,N-1\}\), allowing repetitions, and form
\(A=\sqrt{N/m}(F_N)_{(r_1,\ldots,r_m),:}\). Put \[
\delta_s(A)=\sup_{\substack{x\in\mathbb C^N\setminus\{0\}\\|\operatorname{supp}(x)|\leq s}}
\left|\frac{\|Ax\|_2^2}{\|x\|_2^2}-1\right|,
\qquad
m_*(N,s)=\min\{m\geq1:\Pr[\delta_s(A)\leq1/3]\geq2/3\}.
\] Determine the order of \(m_*(N,s)\), within absolute multiplicative
constants, uniformly for integers \(N\geq2\) and \(2\leq s\leq N\). In
particular, identify the logarithmic factors that are necessary rather
than artifacts of current upper bounds.

The problem fixes the group to cyclic Fourier, the sampling to
independent rows with replacement, and the success and distortion
thresholds to constants. Results for Walsh--Hadamard matrices do not
automatically give the same sharp answer for every cyclic \(N\). This
distinction matters because lower bounds can depend on subgroup
structure.

Fast multiplication by Fourier sketches is valuable in large
least-squares and sparse recovery computations. The source asks how many
Fourier rows suffice for a fixed restricted-isometry tolerance.
Haviv--Regev's upper bound is \(O(s\log^2(s)\log N)\) at fixed
tolerance; known lower bounds leave a logarithmic gap in general.

\subsection{References}\label{references}

\begin{enumerate}
\def\labelenumi{\arabic{enumi}.}
\tightlist
\item
  A. S. Bandeira, \emph{Ten Lectures and Forty-Two Open Problems in the
  Mathematics of Data Science}, Open Problem 6.1, including its explicit
  sampling-with-replacement Fourier model.
  \href{https://people.math.ethz.ch/~abandeira/TenLecturesFortyTwoProblems.pdf}{Author's
  notes}.
\item
  I. Haviv and O. Regev, \emph{The Restricted Isometry Property of
  Subsampled Fourier Matrices}, SODA 2016; Geometric and Functional
  Analysis 27 (2017), pp.~119--142. Main restricted-isometry theorem.
  \href{https://arxiv.org/abs/1507.01768}{Paper}.
\item
  A. S. Bandeira, M. E. Lewis, and D. G. Mixon, \emph{Discrete
  uncertainty principles and sparse signal processing}, Journal of
  Fourier Analysis and Applications 24 (2018), pp.~935--956. Theorem 16
  gives a Fourier sampling lower bound under stated divisibility
  hypotheses. \href{https://arxiv.org/abs/1504.01014}{Paper}.
\end{enumerate}

\subsection{Status check --- 2026-09-08}\label{status-check-2026-09-08}

Searched ``subsampled Fourier RIP sharp logarithmic 2025 2026'',
``cyclic Fourier RIP lower bound'', and the exact upper-bound title;
checked current arXiv records. No uniform sharp-order resolution was
found. Bounds for partial circulant Gaussian operators and Boolean Walsh
transforms concern different random matrices.
\end{document}
